\chapter{Examples}
\label{chap:examples}

\section{Learning examples}

The \filepath{crates/sitas/examples/} directory contains executable
documentation. Small examples isolate one mechanism before larger examples
combine them:

\begin{itemize}
    \item basic executor futures and timers;
    \item readiness-driven TCP helpers;
    \item sharded executor startup and submission;
    \item shard-local values;
    \item CPU placement reporting;
    \item observability snapshots;
    \item standard-file, mailbox-transfer, and \texttt{io\_uring}
          index-building demos.
\end{itemize}

Each example answers a runtime question: what wakes a future, which shard polls
it, what a snapshot records, or how cancellation behaves. Diagnostic output is
intentionally detailed.

\section{Suggested reading path}

The examples are easiest to learn in layers:

\begin{enumerate}
    \item Start with \filepath{executor_sleep.rs},
          \filepath{executor_timeout.rs}, and \filepath{executor_race.rs}.
          These show futures, timers, and combinators before sharding enters
          the picture.
    \item Move to \filepath{async_readable.rs}, \filepath{async_write.rs}, and
          \filepath{async_tcp_pair.rs}. These show readiness and non-blocking
          file descriptor behavior.
    \item Read \filepath{async_tcp_server.rs},
          \filepath{async_tcp_scoped_server.rs}, and
          \filepath{async_tcp_scheduling_groups.rs}. These show handler
          spawning, scoped shutdown, and scheduling groups.
    \item Then read \filepath{sharded_executor.rs},
          \filepath{sharded_submit.rs}, and \filepath{sharded_map_reduce.rs}.
          These show explicit placement of async work on shards.
    \item Finally read \filepath{shard_local.rs} and
          \filepath{shard_local_worker_observability.rs}. These show local
          state access and owned runtime snapshots.
\end{enumerate}

Run an example, then inspect the modules it uses. Most print counters or
snapshots that connect visible behavior to runtime state.

\section{Observability examples}

\filepath{sharded_observability.rs} samples task state while named shard tasks
sleep and wake. The interesting columns are:

\begin{description}
    \item[\texttt{status}] the coarse task lifecycle state;
    \item[\texttt{age\_ms}] how long the task has existed;
    \item[\texttt{state\_ms}] how long it has been in the current state;
    \item[\texttt{polls}] how many times the future has been polled;
    \item[\texttt{wait}] the last known wait interest.
\end{description}

\filepath{async_tcp_scheduling_groups.rs} shows a different form of
observability. It starts two TCP accept loops, parks accepted handlers, takes a
snapshot, and prints which scheduling group each live task belongs to. The
snapshot is owned data, so any tool can render it.

\filepath{scheduling\_group\_demo.rs} and
\filepath{sharded\_scheduling\_groups.rs} expose group-level progress as well:
charged poll count, charged poll time, average poll duration, and poll-time
share. These fields make long-running cooperative work easier to inspect
without adding a tracing subsystem or global scheduler.

\section{Sharded index build}

The sharded index example uses two files: a fixed-size record data file and a
sorted index of offsets into that data file. Each shard scans one partition,
sorts its partition-local index entries, writes a materialized run file, and
then participates in merge rounds until one sorted index remains.

The example demonstrates several runtime APIs in combination:

\begin{itemize}
    \item \texttt{ShardLocal<ShardProgress>} tracks per-shard phase,
          record count, and byte counters. Each shard task calls
          \texttt{with\_current\_result} to update only its own slot.
          No shared mutex protects the progress counters.
    \item \texttt{runtime.map\_all} dispatches partition work to all
          shards with a single call, replacing manual handle collection.
    \item Chunked record reads (256 records per read) replace one-syscall-
          per-record scanning.
    \item \texttt{RunningStatistics} collects per-shard scan durations and
          prints mean, standard deviation, CV, min, and max across shards.
    \item Explicit merge rounds submit pairs of sorted run files onto shards,
          stream-merge them into new run files, and carry odd runs forward.
    \item Wall-clock timing is reported at the end:
          \texttt{partition=34.1ms merge=82.2ms total=175.7ms}.
\end{itemize}

The \texttt{io\_uring} variant performs partition reads, run writes, and merge
run reads/writes through the executor-backed Linux completion path. It uses the
same \texttt{ShardLocal} and \texttt{map\_all} shape as the std variant. Final
verification still uses ordinary file I/O because it checks the externally
visible result rather than the runtime path that produced it.

The index examples split work by shard, keep partition state local, materialize
intermediate runs, merge owned outputs, and collect per-shard timing data. Read
them after the smaller executor and sharding examples.

The index build is a good larger example because it is naturally partitionable.
Each shard sorts a portion of the data without sharing a mutable index
structure. The merge phase makes coordination cost explicit by exchanging owned
run files and summaries rather than borrowing another shard's working set.

\subsection{Mailbox-transfer index build}

\filepath{sharded\_index\_mailbox.rs} uses a different assembly path from
\filepath{sharded\_index\_build.rs}. Scan tasks still run on executor shards,
but they send owned \texttt{IndexMessage} batches through
mailbox senders instead of using intermediate run files as the cross-shard exchange medium.

The demo separates logical ownership from physical shard placement:

\begin{itemize}
    \item \texttt{WorkUnitRouter<IndexWorkUnit>} maps each record key to a
          logical assembler work unit.
    \item \texttt{WorkUnitMailboxSet<IndexWorkUnit, IndexMessage>} maps each
          assembler name to the shard that owns its receiver.
    \item The default run starts one more assembler than shards, so one shard
          owns two logical assembler work units. This makes non-uniform
          placement visible in the printed placement table.
    \item Each assembler receives owned batches, sorts its own entries, writes
          one final partition file, and returns owned metadata.
\end{itemize}

The communication paths are mixed. Scanner tasks send \texttt{Vec<IndexEntry>} batches to assembler work units through mailboxes.
Assemblers do not send mailbox messages to the coordinator. Instead, each
assembler returns small owned \texttt{PartitionRun} metadata through its task
handle, and the coordinator reads the assembler run files during the final
k-way merge.

The final index is a file, while intermediate cross-shard exchange uses typed,
owned messages rather than shared mutable state.

\subsubsection{Performance aside: mailbox transfer vs file-backed exchange}

The two index variants produce identical results but differ in how intermediate
data crosses shard boundaries. The build example writes per-shard sorted run
files and performs multi-round pairwise merge I/O. The mailbox example routes
entries by key through owned messages to assembler work units, which sort and
write pre-partitioned files before a single k-way merge.

Measured on macOS (Apple Silicon, 8 CPUs), release mode, deterministic input:

\begin{center}
\rowcolors{2}{sitasCodeBackground}{white}
\begin{tabular}{r r r r r r}
\hline
\textbf{Records} & \textbf{Shards} & \textbf{Build total} & \textbf{Mailbox total}
    & \textbf{Build merge} & \textbf{Mailbox scan+xfer} \\
\hline
1\,M & 1 & 6.1\,s  & 8.9\,s  & \(\approx\)0 & 20\,ms \\
1\,M & 2 & 8.6\,s  & 8.9\,s  & 2.9\,s  & 13\,ms \\
1\,M & 4 & 10.2\,s & 8.2\,s  & 5.3\,s  & 9\,ms  \\
1\,M & 8 & 13.3\,s & 9.3\,s  & 6.9\,s  & 12\,ms \\
5\,M & 2 & 47.2\,s & 46.3\,s & 15.0\,s & 63\,ms \\
5\,M & 4 & 57.4\,s & 45.4\,s & 27.4\,s & 39\,ms \\
5\,M & 8 & 69.9\,s & 51.0\,s & 34.9\,s & 43\,ms \\
\hline
\end{tabular}
\end{center}

In this measurement, the build variant slows as shard count increases because
more run files require more full-dataset merge rounds. At 5\,M records, moving
from 2 to 8 shards increased its total by 48\%. The mailbox variant stayed
closer to flat because it replaced intermediate file exchange with in-memory
transfer and one final k-way merge.

Both totals are dominated by unbuffered I/O overhead. \texttt{write\_index\_entry}
performs two 8-byte \texttt{write\_all} calls per entry. At 5\,M records that
is 10\,M write syscalls for the index alone. Data file creation and
verification suffer the same pattern. The mailbox's reported work phases
(scan+transfer in tens of milliseconds, assembler join in microseconds)
reflect concurrent shard execution where scanning, mailbox transfer, assembler
sort, and partition writes overlap on executor threads.

For this workload, the build variant wins at one shard because it has no merge
phase. At four or more shards, the mailbox variant avoids the pairwise merge
I/O. These measurements are illustrative rather than general performance
claims; both variants use unbuffered per-entry writes.

\section{Representative commands}

\begin{lstlisting}[language=bash,caption={Local validation}]
cargo fmt --check
cargo clippy --all-targets -- -D warnings
cargo test
cargo doc --no-deps
cargo run --example sharded_index_mailbox -- --records 512 --shards 4
\end{lstlisting}

\begin{lstlisting}[language=bash,caption={Linux validation through Docker}]
SITAS_DOCKER_IO_URING=1 tools/linux-docker.sh cargo test
SITAS_DOCKER_IO_URING=1 tools/linux-docker.sh cargo run --example sharded_index_build_uring -- --records 512 --shards 4 --seed 0xbeef --cleanup
\end{lstlisting}
